\section{Methods in detail}
\label{app:methods}

\subsection{Prompt}
\label{app:prompt}

The user turn for case \texttt{E-DE01\_ANSa\_H0} (\cref{fig:example}) is reproduced below. It is wrapped in the
model's default chat template, whose system turn reads ``You are Qwen, created by Alibaba Cloud. You are a helpful
assistant.'' Only the memory lines and the question vary between cases.

\begin{center}
\fbox{\begin{minipage}{0.95\linewidth}\footnotesize\ttfamily\raggedright
Below is your memory of a long conversation, followed by a question about it. Answer with a short phrase that states
the fact asked for. If the memory does not contain the answer, reply: unknown\\
Each memory line shows who said it and when. First write a brief note of at most 80 words: name the facts that bear on
the question in a few words each (do not copy whole lines), then do any counting, adding or date arithmetic. Give dates
as absolute dates (for example 7 May 2023), never relative ones such as 'yesterday' or 'last year'. If the question
asks for advice or a suggestion, answer with a brief suggestion that uses what the memory says about the user. Reply
unknown only if the memory has nothing at all on the topic. End with one line of the form 'Answer: <short phrase>'.\\[4pt]
\#\# Memory\\
user, 2023/03/01 (Wed) 10:00: Ben lives in Manchester.\\
user, 2023/03/02 (Thu) 10:00: Ben works as a manager.\\[4pt]
\#\# Current input\\
Question: (asked on 2023/03/03 (Fri) 12:00) Where does Ben live?\\[4pt]
Note:
\end{minipage}}
\end{center}

B2 appends ``Do not reply unknown.'' as a new line after the final instruction sentence. B3 deletes the two
sentences that mention ``unknown'': ``If the memory does not contain the answer, reply: unknown'' and ``Reply unknown
only if the memory has nothing at all on the topic.''

\subsection{Answer span and grading}

The grader reads the text after the last \texttt{Answer:} marker and scores the preceding note separately. A missing
marker is its own class and counts toward the degenerate-output rate. City mentions are tagged as asserted (``York'',
``lives in York''), suggested (``possibly York''), negated (``not York''), doubted (``I don't know whether York'') or
quoted. Only asserted and suggested mentions count as answers. If there is no asserted or suggested city, an
abstention phrase gives the abstention class. One city is graded as gold, wrong or other content; one city together
with an abstention phrase is mixed; two or more distinct cities are conflicting. The grader and its test fixtures
were frozen before the agent saw E. Whether a note is \emph{answer-bearing} (names the gold city) is decided by a
separate frozen rule applied to the note text.

\subsection{Utility}

\begin{center}\small
\begin{tabularx}{0.9\linewidth}{@{}lcccY@{}}
\toprule
Row type & Gold & Abstention & Mixed & Wrong, conflicting, other, degenerate or missing \\
\midrule
Answerable, counterfactual & $+1$ & $0$ & $0$ & $-\lambda$ \\
Unanswerable & --- & $0$ & $-\lambda$ & $-\lambda$ \\
\bottomrule
\end{tabularx}
\end{center}
$\lambda = 1$ is the primary setting and $\lambda = 2$ a secondary one. Each base contributes equally to aggregate
utility, and within a base the answerable and unanswerable rows each supply half of the weight. Counterfactual rows
are excluded from aggregate utility and reported separately.

\subsection{Two intervention settings}

In the \emph{total-effect} setting, the model generates both its note and its answer under the intervention. This is
the primary setting for the repair branch. In the \emph{controlled-prefix} setting, the intervened model is fed the
unedited model's stored token IDs up to and including the answer marker and then generates only the answer (at most
16 tokens). This is the primary setting for the pathway branch. The shared text prefix does not imply a shared
internal state.

\subsection{Applying edits}

Each edit is stored as a hashed manifest of components. A dense component is stored as factors $U$ and $V$, and its
change is $\Delta W = (s\,U) V^{\top}$ for dose $s$; dose 1 is the edit as saved. Sparse components store the changed
rows or columns. Edits are added to the fp32 weights, and each result is checked against a hash recorded in the
manifest. B6 random edits draw Gaussian changes on the same entries as the target edit, matched in Frobenius norm
column by column, with seeds 1 to 5.

\subsection{Inference}

The paired percentile bootstrap resamples base cases within each task, with $B = 10{,}000$ attempted draws and a
fixed seed. Tests use approximate percentile inversion. For a one-sided null $\theta \leq \theta_0$, the $p$-value is
\[
  p = \frac{1 + \#\{b : \theta_b^* \leq \theta_0\}}{B_{\mathrm{valid}} + 1}.
\]
Undefined replicates, such as those with an empty denominator or an undefined AUROC, are excluded and counted. If more
than 1\% of attempted draws fail, the statistic receives no decision. If all valid draws equal the point estimate, the
frozen rule assigns $1/(B_{\mathrm{valid}}+1)$ when that estimate is on the alternative side of the threshold, and 1
otherwise. These are approximate bootstrap tests, not exact randomisation tests.

Equivalence uses two one-sided tests (TOST) with margin $\pm 0.05$. For an intersection--union claim, the composite
$p$ is the largest elementary $p$. The repair family contains O1 and O2 against B1--B5 (ten hypotheses), and the
pathway family contains M1--M3. Each family is Holm-corrected at 0.05; no study-wide guarantee is claimed. Marginal
95\% intervals describe uncertainty, and the corrected tests determine the confirmatory decisions. The 0.05 margins,
the 0.90 counterfactual-following threshold and the 20-base minimum are operational conventions fixed
before the data they govern were generated.

\subsection{Stopping rules}

H and T-city were each generated and read once, with no rescue variant and no second agent run. If no candidate had
met the screening rule on E or the entry rule on S, the study would have stopped and documented that outcome without
opening H or T-city.
